\begin{textsample}{NEG Dim 3 – 2010\_c6 – Score -7.00 – 2010\_c6/enade\_2010\_odonto\_d1b7c708cc.txt}  \label{ex:f3_neg_013}
CAPÍTULO 3 ANÁLISE TÉCNICA DA PROVA Este capítulo tem por objetivo apresentar o desempenho dos estudantes de Odontologia no ENADE/2010.
Para isso, foram calculadas as estatísticas básicas da prova em seu todo, bem como as estatísticas dos componentes relacionadas à Formação Geral e ao de Conhecimento Específico da Área.
Nas tabelas, são apresentados o tamanho da população inscrita e de presentes, e as seguintes estatísticas das notas7 : média do desempenho na prova, erro padrão da média, desvio padrão, nota mínima, mediana, nota máxima, coeficiente de assimetria e de variação.
As estatísticas apresentadas neste capítulo contemplam o total de estudantes e, separadamente, os Ingressantes e os Concluintes.
Foram calculadas tendo -se em vista as seguintes agregações : ( a ) as Grandes Regiões e o país como um todo ; ( b ) a Categoria Administrativa ; e ( c ) a Organização Acadêmica.
Em relação aos gráficos de distribuição de notas, o intervalo considerado foi de 10 unidades, aberto à esquerda e fechado à direita, com exceção do primeiro intervalo, [ 0 ; 10 ], fechado em ambos os extremos.
Para os gráficos de distribuição das notas das questões discursivas, foram consideradas mais duas categorias : questão em branco e nota zero. 3.1 ESTATÍSTICAS BÁSICAS DA PROVA 3.1.1 Estatísticas Básicas Gerais A Tabela 3.1 apresenta as estatísticas básicas da prova por grupo de estudantes.
A população total de estudantes regularmente inscritos8 para a prova foi 20.099.
Destes, 17.587 estiveram presentes, sendo 12,5 \% o índice de não comparecimento.
A abstenção foi maior entre Ingressantes ( 17,6 \% ) do que entre Concluintes ( 2,8 \% ).
A nota média de toda a prova foi 42,5, sendo que os Ingressantes obtiveram média mais baixa, 35,5, que os Concluintes, cuja média foi 53,8.
O desvio padrão foi 14,7, sendo o do grupo de Ingressantes menor ( 11,4 ) do que o do grupo dos Concluintes ( 12,1 ) ; os coeficientes de variação ( 32,1 \% no grupo dos Ingressantes e 22,5 \% no dos Concluintes ) indicam menor dispersão relativa nas notas de Concluintes.
A nota máxima foi 87,7, obtida por Concluintes, ao passo que a maior nota obtida por Ingressantes foi 71,1. 7 Algumas dessas estatísticas estão definidas no Capítulo 1. 8 As estatísticas das notas apresentadas consideram apenas os estudantes regularmente inscritos e com resultados validados por ausência de erro administrativo.
Uma comparação entre o desempenho de Ingressantes e Concluintes também pode ser realizada a partir dos dados apresentados no Gráfico 3.1, com a distribuição de notas finais.
Destaca -se que as notas mais elevadas, considerando -se o intervalo posterior à nota 40,0, foram atingidas principalmente pelos Concluintes, cuja concentração das notas está à direita dos Ingressantes.
Já nas faixas de notas mais baixas, inverteu -se o quadro identificado nas faixas de notas mais altas.
Essa característica é corroborada pelo valor da nota mediana 35,6 no grupo dos Ingressantes, abaixo da encontrada no grupo dos Concluintes, 54,8, transparecendo uma superioridade de desempenho dos Concluintes em relação aos Ingressantes na prova em tela ( Tabela 3.1 ).
Os coeficientes de assimetria indicam uma assimetria negativa ( -0,5 ) para os Concluintes, e comportamento simétrico ( 0,0 ) para Ingressantes.
O Gráfico 3.2 apresenta informações referentes à nota final de Ingressantes e Concluintes, comparando -se os resultados de acordo com as Grandes Regiões do país, a Categoria Administrativa e a Organização Acadêmica.
Considerando -se as notas médias dos alunos em cada Grande Região, observa -se que, em relação às notas aferidas em âmbito nacional ( média 35,5 para Ingressantes e 53,8 para Concluintes ), existe diferença estatisticamente significativa ao nível de 95 \% entre a menor média, obtida na região Norte ( 32,4 ) e a maior, na região Sul ( 37,4 ), no grupo de Ingressantes.
Entre os Concluintes, também existe diferença estatisticamente significativa das médias ao nível de 95 \%.
A maior média de notas foi obtida pela região Sul ( 56,4 ), ao passo que a menor foi identificada na região Norte ( 49,3 ).
As demais médias de notas foram 53,2 para a região Sudeste, 53,5 para a região Centro-Oeste e 54,9 para a região Nordeste.
Destaca -se a região Sul com médias superiores às nacionais, tanto para o grupo de Ingressantes quanto para o de Concluintes.
Levando -se em conta os grupos dos estudantes em Categorias Administrativas, observa -se que a pontuação média mais elevada entre os Concluintes foi encontrada em instituições públicas ( 58,8 ), com desempenho acima da média nacional.
Já os Concluintes provenientes de instituições privadas obtiveram média 51,1, abaixo da média nacional.
Quanto aos Ingressantes, as notas mais elevadas, também acima da média nacional para esse grupo de estudantes, ocorreram nas instituições públicas ( 41,3 ).
Os Ingressantes vinculados a instituições privadas obtiveram a nota média 32,9, abaixo da média nacional.
Tendo como foco as notas médias dos estudantes agrupados por Organização Acadêmica, distribuídas entre Universidades, Centros Universitários e Faculdades, para os Concluintes vinculados a Universidades, a média 55,8 ficou acima da média nacional.
Pode -se verificar a mesma realidade entre os Ingressantes de Universidades, cuja nota média foi 37,3, também acima da nacional para esse grupo.
Considerando -se as demais Organizações Acadêmicas, observa -se que nos Centros Universitários e Faculdades, o desempenho de Ingressantes e de Concluintes ficou abaixo da média nacional, com notas médias 49,4 e 48,9, respectivamente.
Ingressantes das mesmas Organizações Acadêmicas obtiveram médias 32,7 e 31,9, respectivamente. 3.1.2 Estatísticas Básicas no Componente de Formação Geral A Tabela 3.2 apresenta as estatísticas básicas em relação ao componente da prova que avalia a Formação Geral dos estudantes.
Os Concluintes obtiveram desempenho médio 49,2, que foi superior ao dos Ingressantes, 43,9.
Quanto à variabilidade das notas nos dois grupos, o desvio padrão das notas obtidas foi 18,0 para Ingressantes e 16,8 para Concluintes.
Os dois grupos alcançaram notas máximas semelhantes, sendo 96,0 a maior nota obtida por pelo menos um Ingressante, e 95,0 a maior nota entre os Concluintes.
O Gráfico 3.3 propicia a comparação entre Concluintes e Ingressantes relativamente ao desempenho no componente de Formação Geral.
Nesse componente, os Concluintes apresentaram desempenho superior.
Dos Ingressantes, 19,3 \% obtiveram notas na faixa entre 40 e 50 pontos, e percentual ligeiramente superior na faixa seguinte, de 50 a 60 pontos, a modal : 19,5 \%.
Nos intervalos de notas posteriores, mais elevados, verifica -se queda gradual de ocorrência de Ingressantes.
Entre os Concluintes, houve também um percentual maior de estudantes no intervalo entre 50 e 60 pontos, com 23,6 \%.
Também observou -se queda gradual no desempenho nos intervalos de notas mais elevados, com 18,0 \% no intervalo entre 60 e 70 pontos e 7,8 \% atingido a faixa de 70 a 80.
Apenas 1,4 \% atingiram notas entre 80 e 90 pontos.
Na menor faixa de notas, que vai de zero a 10 pontos, 2,1 \% dos Concluintes ficaram situados.
Acima de 40 pontos, os Concluintes apresentam uma maior representatividade do que os Ingressantes, o inverso ocorrendo para as faixas de nota mais baixas.
O coeficiente de assimetria foi negativo tanto para Concluintes ( -0,4 ) quanto para Ingressantes ( -0,2 ).
No Gráfico 3.4 são apresentadas as informações referentes ao desempenho dos Ingressantes e dos Concluintes no componente de Formação Geral, em diferentes agregações como Grande Região do país, Categoria Administrativa e Organização Acadêmica.
Considerando -se as notas médias dos estudantes segundo cada Grande Região, observa -se que no grupo de Ingressantes existe diferença estatisticamente significativa ao nível de 95 \% entre as médias na região Nordeste ( 47,9, a mais alta ) e Norte ( 41,7, a mais baixa ).
No grupo de Concluintes, as notas mais elevadas ocorreram nas regiões Sul ( 51,7 ) e Nordeste ( 51,2 ).
A menor nota entre Concluintes foi 46,7, correspondente também aos estudantes da região Norte, repetindo o padrão dos Ingressantes.
Essa região ficou abaixo da média nacional ( 43,9 para Ingressantes e 49,2 para Concluintes ) no componente de Formação Geral.
O mesmo foi verificado na região Sudeste, que obteve nota média 42,3 para Ingressantes e 47,8 para Concluintes.
Tendo como foco as Categorias Administrativas, as instituições públicas obtiveram médias mais elevadas do que as privadas : para Concluintes, notas médias de 54,7 em cursos de IES públicas e de 46,1 nos das privadas.
Para os Ingressantes, as notas foram de, respectivamente, 53,1 e 39,7.
Em comparação à média nacional, as notas de Ingressantes e de Concluintes vinculados às instituições públicas foram superiores, ao passo que, nas instituições privadas, foram inferiores.
Considerando -se o tipo de Organização Acadêmica, nota -se diferença estatisticamente significativa ao nível de 95 \% entre a melhor e a menor média.
Nas Universidades, houve desempenho melhor do que em Centros Universitários e Faculdades.
Neste primeiro tipo de Organização Acadêmica, Ingressantes tiveram nota média 46,2, enquanto Concluintes, 51,0.
Nos Centros Universitários e Faculdades, os valores foram mais baixos : os Ingressantes tiveram média 38,7 nos Centros Universitários e 40,4 nas Faculdades ; entre Concluintes, a média foi 44,1 para Centros Universitários e 45,4 para as Faculdades. 3.1.3 Estatísticas Básicas do Componente de Conhecimento Específico A Tabela 3.3 apresenta as estatísticas básicas referentes à parte do componente de Conhecimento Específico da Área de Odontologia.
A média total de desempenho foi 41,3.
Os Concluintes apresentaram desempenho médio ( 55,3 ) superior ao dos Ingressantes ( 32,7 ) ; as medidas de dispersão – desvio padrão e amplitude geral das notas – foram menores no grupo dos Ingressantes do que no dos Concluintes.
Os coeficientes de variação, 36,4 \% no grupo dos Ingressantes e 23,0 \% no dos Concluintes, indicam dispersão relativa menor nas notas desses últimos.
As notas máximas obtidas por Concluintes e Ingressantes foram iguais a 90,4 e 74,1, respectivamente, notas inferiores àquelas encontradas no componente de Formação Geral ( 96,0 e 95,0, respectivamente ).
As medianas das distribuições das notas dos Ingressantes e dos Concluintes foram 33,2 e 56,3, respectivamente, enfatizando a superioridade dos Concluintes nesse componente da prova.
Assim como os Gráficos 3.1 e 3.3, o Gráfico 3.5, apresentado a seguir, proporciona uma comparação do desempenho entre Concluintes e Ingressantes, em relação ao componente de Conhecimento Específico.
Dentre as três comparações feitas, observa -se que, embora a distribuição das notas seja semelhante à apresentada na prova como um todo, neste componente a diferença de desempenho entre Concluintes e Ingressantes é mais acentuada, destacando -se o desempenho mais elevado dos primeiros.
Claramente a distribuição das notas dos Ingressantes, mais concentradas à esquerda em relação aos Concluintes, evidencia o percentual mais alto de notas nas faixas mais baixas para este grupo.
A maior concentração de Concluintes está nos intervalos compreendidos entre ( 50 a 60 ] pontos, ao passo que a de Ingressantes está entre ( 20 a 30 ] pontos.
Tais como os Gráficos 3.2 e 3.4, o Gráfico 3.6 apresenta uma comparação dos resultados em relação à Grande Região do país, à Categoria Administrativa e à Organização Acadêmica, agora levando em conta o desempenho dos Ingressantes e dos Concluintes no componente de Conhecimento Específico da prova.
Ao se observar as notas médias dos estudantes, segundo cada Grande Região do país, verifica -se que, entre os Concluintes, a mais elevada foi encontrada na região Sul ( 57,9 ), maior que a média nacional ( 55,3 ).
Em relação aos Ingressantes, também a região Sul apresenta a média mais elevada ( 34,8 ), seguida da região Sudeste, com média igual a 32,9.
As médias mais baixas entre os Concluintes ( 50,2 ) e entre os Ingressantes ( 29,3 ) foram verificadas na região Norte.
Quanto à Categoria Administrativa, observa -se comportamento semelhante àquele da parte relativa à Formação Geral e à prova como um todo, ou seja, as médias mais elevadas tanto entre os Concluintes quanto entre os Ingressantes foram encontradas em instituições de ensino públicas.
Nesse tipo de IES, a nota média de Concluintes foi 60,1, enquanto a de Ingressantes, 37,3.
Tais médias estão situadas acima da média nacional.
Nas instituições privadas, as médias foram, mais uma vez, inferiores em relação às públicas.
Concluintes obtiveram média 52,7, e Ingressantes, 30,6.
Em relação à Organização Acadêmica, o desempenho dos Concluintes foi superior ao dos Ingressantes nas três categorias.
Nas Universidades, os Concluintes obtiveram a média mais elevada, igual a 57,3.
Mais abaixo estiveram os Centros Universitários ( 51,2 ) e as Faculdades ( 50,0 ).
Os Ingressantes de Universidades também tiveram nota mais elevada ( 34,3 ) que a nacional.
Os estudantes de Centros Universitários obtiveram média 30,6, próxima aos de Faculdades, com 29,0. e as Faculdades ( 50,0 ).
Os Ingressantes de Universidades também tiveram nota mais elevada ( 34,3 ) que a nacional.
Os estudantes de Centros Universitários obtiveram média 30,6, próxima aos de Faculdades, com 29,0. 3.2 ANÁLISE DAS QUESTÕES OBJETIVAS 3.2.1 Componente de Formação Geral A Tabela 3.4 apresenta as estatísticas básicas relativas às oito questões objetivas do componente da prova que abrange a Formação Geral dos estudantes.
Como pode ser observado, os estudantes Concluintes obtiveram desempenho superior, com nota média 51,6, em relação aos alunos Ingressantes, cuja nota média foi 46,1.
As medidas de dispersão nos dois grupos de alunos foram similares, como indicam os desvios-padrão, assim como a amplitude total das notas, sendo que a nota mínima foi zero e a máxima, 100,0, para os dois grupos de alunos avaliados.
Os coeficientes de variação, 41,4 \% no grupo dos Ingressantes e 35,1 \% no grupo dos Concluintes, indicam menor dispersão relativa nas notas desses últimos.
A mediana foi igual a 50,0 para os dois grupos.
O coeficiente de assimetria do total de estudantes revelou assimetria negativa ( -0,2 ).
A Tabela 3.5 apresenta as classificações das questões objetivas do componente de Formação Geral, segundo o índice de facilidade.
As questões objetivas da prova foram assim avaliadas : das oito questões, nenhuma teve o índice de facilidade classificado como muito fácil.
Quatro questões foram tidas como fáceis, por terem índice de acertos situado na faixa entre 0,61 e 0,85 ( de 61,0 \% a 85,0 \% de acertos ).
Uma questão foi considerada de dificuldade média, situando -se no intervalo entre 0,41 e 0,60 do índice de facilidade, ou seja, houve entre 41,0 \% e 60,0 \% de acertos.
Uma questão foi classificada como difícil, situando -se no intervalo entre 0,16 e 0,40.
Por fim, duas questões apresentaram menos de 0,15 de acertos, razão pela qual foram classificadas como muito difíceis.
A Tabela 3.6 descreve os resultados das análises das questões objetivas relativas à Formação Geral, segundo o poder de discriminação, utilizando -se para tal, do índice de discriminação ponto bisserial.
Cinco das oito questões apresentaram índices acima de 0,40 e, assim, foram classificadas como muito boas nesta característica para esse grupo de alunos ; uma questão teve um bom índice de discriminação, entre 0,30 e 0,39.
Foram classificadas como de média discriminação duas das questões.
Nenhuma delas teve fraco índice de discriminação.
Na sequência, na Tabela 3.7, estão especificadas, por questão, os resultados da análise relativa aos índices de facilidade e de discriminação de cada questão no componente relativo à Formação Geral da prova de Odontologia.
O índice de facilidade variou de 0,11 a 0,80, e o de discriminação, de 0,26 a 0,55.
As questões com índices de discriminação muito bom, de números 1, 3, 4, 5 e 8, figuraram entre as mais fáceis desse conjunto : a questão 1 com dificuldade média e as demais classificadas de fácil.
A questão 8 foi a que apresentou maior poder discriminatório, com índice 0,55, e foi também a mais fácil, com uma proporção de 80,0 \% de acertos.
A questão de número 6 apresentou índice de facilidade 0,11, ou seja, um quantitativo pequeno de estudantes conseguiu resolvê -la, dentro do universo de A Tabela 3.6 descreve os resultados das análises das questões objetivas relativas à Formação Geral, segundo o poder de discriminação, utilizando -se para tal, do índice de discriminação ponto bisserial.
Cinco das oito questões apresentaram índices acima de 0,40 e, assim, foram classificadas como muito boas nesta característica para esse grupo de alunos ; uma questão teve um bom índice de discriminação, entre 0,30 e 0,39.
Foram classificadas como de média discriminação duas das questões.
Nenhuma delas teve fraco índice de discriminação.
Na sequência, na Tabela 3.7, estão especificadas, por questão, os resultados da análise relativa aos índices de facilidade e de discriminação de cada questão no componente relativo à Formação Geral da prova de Odontologia.
O índice de facilidade variou de 0,11 a 0,80, e o de discriminação, de 0,26 a 0,55.
As questões com índices de discriminação muito bom, de números 1, 3, 4, 5 e 8, figuraram entre as mais fáceis desse conjunto : a questão 1 com dificuldade média e as demais classificadas de fácil.
A questão 8 foi a que apresentou maior poder discriminatório, com índice 0,55, e foi também a mais fácil, com uma proporção de 80,0 \% de acertos.
A questão de número 6 apresentou índice de facilidade 0,11, ou seja, um quantitativo pequeno de estudantes conseguiu resolvê -la, dentro do universo de participantes.
A questão 7 obteve índice de discriminação médio, com 0,26, e a de número 2 obteve índice 0,37, considerado bom.
O Gráfico 3.7 analisa o comportamento da questão de número 8 de Formação Geral.
Trata -se da questão mais fácil e com o maior índice de discriminação desse componente da prova.
No gráfico a seguir, cada uma das cinco curvas representa o percentual de respostas em determinada alternativa da questão, em função da nota dos estudantes nesta parte da prova ( Formação Geral/Múltipla Escolha ) antes de possíveis eliminações pelo critério do ponto bisserial.
A curva em rosa corresponde à alternativa B, a correta para este item.
Assim, observa -se que entre os estudantes com notas mais baixas, desempenho mais fraco nesse componente do exame, a situação mais frequente foi a escolha de uma das alternativas incorretas : a alternativa C ( em verde ) ou E ( em roxo ).
Na medida em que a nota aumenta, indicando desempenho melhor nesta parte da prova, aumenta concomitantemente a proporção de estudantes que selecionaram a alternativa correta B, atingindo 100 \% para as notas acima de 80,0.
Essa análise permite verificar como a questão discriminou entre os grupos de desempenho, justificando o alto índice obtido.
Os gráficos relativos às demais questões de Formação Geral constam do Anexo I. 3.2.2 Componente de Conhecimento Específico A Tabela 3.8 apresenta as estatísticas básicas em relação às questões objetivas do componente de Conhecimento Específico da prova.
Nesse conjunto de questões, os estudantes Concluintes ( 57,9 ) obtiveram nota média superior, aos Ingressantes ( 37,2 ).
Quanto à variabilidade das notas, o desvio padrão no grupo dos Ingressantes e de Concluintes foi igual, 13,3.
Os coeficientes de variação foram 35,8 \% e 23,0 \%, respectivamente, indicando menor dispersão relativa nas notas dos Concluintes.
Com relação à mediana, no grupo dos Ingressantes, esta foi igual a 39,1, estando acima da média desse grupo, enquanto no grupo dos Concluintes, a mediana foi igual a 60,9, também superando um pouco a média desse grupo.
Ingressantes tiveram comportamento simétrico ( 0,0 ) e o coeficiente de assimetria foi negativo ( -0,5 ) para Concluintes.
A Tabela 3.9 apresenta as classificações das questões objetivas da parte da prova relativa ao componente de Conhecimento Específico, segundo o índice de facilidade.
A partir dos índices obtidos, pode -se concluir que, das questões objetivas da prova, a maior parte foi de média a difícil : das 27 questões, doze foram classificadas como difíceis e duas como muito difíceis e nove delas foram classificadas como médias.
Três questões foram tidas como fáceis e apenas uma como muito fácil.
Na Tabela 3.10 são apresentados os índices de discriminação das questões objetivas do componente de Conhecimento Específico da prova.
Observa -se, segundo os índices estabelecidos, que a capacidade de discriminação das questões objetivas desse componente da prova foi classificada, para nove das 27 questões, como muito boa, enquanto seis delas tiveram índice de discriminação bom.
Uma boa quantidade de questões – 15 em 27 – proporcionou, portanto, índices de discriminação bons ou muito bons.
Oito delas foram classificadas como médias e outras quatro como fracas, sendo doze, por conseguinte, a quantidade de questões nos dois patamares mais baixos de discriminação.
Na Tabela 3.11 estão especificados, por questão, os resultados da análise relativa aos índices de facilidade e de discriminação de cada uma das questões do componente de Conhecimento Específico da prova de Odontologia.
Dentre as questões que alcançaram os maiores índices de discriminação, muito bom, as de números 15, 19, 22, 23, 26, 29, 33, 36 e 37, uma ( a questão 33 ) foi classificada como muito fácil pelo índice de facilidade, três ( questões 23, 36 e 37 ) foram classificadas como fáceis, quatro como de dificuldade média ( questões 19, 22, 26 e 29 ) e a questão 15 como difícil.
Em termos de proporção de acertos, tais questões variaram de 0,33 a 0,76.
Quanto aos índices de discriminação, estiveram entre 0,40 e 0,52. discriminação.
Estas quatro questões ( 14, 17, 24 e 34 ) com índice de discriminação fraco foram eliminadas do cômputo da média final das notas.
Na Tabela 3.11 estão especificados, por questão, os resultados da análise relativa aos índices de facilidade e de discriminação de cada uma das questões do componente de Conhecimento Específico da prova de Odontologia.
Dentre as questões que alcançaram os maiores índices de discriminação, muito bom, as de números 15, 19, 22, 23, 26, 29, 33, 36 e 37, uma ( a questão 33 ) foi classificada como muito fácil pelo índice de facilidade, três ( questões 23, 36 e 37 ) foram classificadas como fáceis, quatro como de dificuldade média ( questões 19, 22, 26 e 29 ) e a questão 15 como difícil.
Em termos de proporção de acertos, tais questões variaram de 0,33 a 0,76.
Quanto aos índices de discriminação, estiveram entre 0,40 e 0,52.
A questão de número 17 foi a mais difícil dentre as 27 questões específicas, com baixo índice de facilidade, 0,10, o que, em termos percentuais, corresponde a 10,0 \% de acertos.
Essa questão apresentou poder discriminatório igualmente baixo, 0,08, o que comprova ter sido esta a mais difícil para os estudantes e foi eliminada do cômputo da média final de notas, como já mencionado.
Destaca -se, também, a questão 14, para a qual o índice de facilidade foi 0,11, o que, em termos percentuais, corresponde a 11,0 \% de estudantes que responderam acertadamente, sendo também uma das mais difíceis da prova.
Obteve índice de discriminação fraco, 0,14 e foi igualmente eliminada.
O Gráfico 3.8 analisa o comportamento da questão 33 do componente de Conhecimento Específico.
Esta foi a questão mais fácil da prova, apresentando índice de facilidade 0,86, ou seja, 86,0 \% dos estudantes assinalaram acertadamente a opção A, correspondente ao gabarito.
Seu índice de discriminação foi igual a 0,42, classificado como muito bom.
Neste gráfico, cada uma das cinco curvas representa o percentual de respostas em determinada alternativa da questão 31, em função da nota dos estudantes nesta parte da prova antes de possíveis eliminações de questões pelo critério do ponto bisserial.
A alternativa correta A, representada no gráfico pela curva em azul marinho, foi escolhida em maiores proporções pelos alunos com desempenho melhor nesta parte da prova, enquanto as alternativas incorretas, também denominadas distratores, foram selecionadas principalmente por aqueles com notas mais baixas, ainda que aqueles com nota zero, na sua quase totalidade tenham deixado em branco ou marcaram mais de uma alternativa, comportamento considerado inválido.
A proporção de alunos que selecionaram a resposta correta A aumenta gradativamente, com a nota dos alunos chegando a atingir 100 \% para notas acima de 70, enquanto a proporção dos que escolhem alternativas incorretas, decai a partir da primeira nota não nula como função da nota nesta parte da prova.
Os gráficos relativos às demais questões do componente de Conhecimento Específico constam do Anexo I. 3.3 ANÁLISE DAS QUESTÕES DISCURSIVAS 3.3.1 Componente de Formação Geral A análise dos resultados de desempenho dos estudantes de Odontologia, nas duas questões discursivas relativas à Formação Geral, encontra -se na Tabela 3.12 e no Gráfico 3.9.
Observa -se, a partir destes, que as notas médias foram mais baixas nesse conjunto do que no conjunto das objetivas.
A média geral dos Ingressantes nas questões objetivas de Formação Geral foi 46,1 ( Tabela 3.4 ), ao passo que nas questões discursivas, a média foi 40,5 para esse grupo.
O mesmo ocorreu entre os Concluintes, que tiveram, em Formação Geral, média 51,6 nas questões objetivas e 45,5 nas questões discursivas.
Cabe destacar que, nas questões discursivas, o desempenho dos Concluintes foi superior ao dos Ingressantes.
Os dois agrupamentos atingiram a mesma nota máxima, 100,0, o que pode sugerir que a superioridade dos Concluintes não tenha sido tão acentuada, em relação aos Ingressantes, nesta parte do exame.
Os coeficientes de assimetria indicam assimetria negativa para Ingressantes e Concluintes.
Verifica -se a semelhança das distribuições das notas desses grupos, no Gráfico 3.9, apresentado a seguir.
Na sequência, os resultados verificados para cada uma das questões discursivas de Formação Geral serão apresentados, estabelecendo -se relações com os conteúdos englobados em cada uma delas.
Os comentários da banca corretora a respeito do observado nas respostas dos estudantes, suas impressões e conclusões serão apresentados junto à análise de cada questão.
Cumpre esclarecer que, tendo em vista que as questões discursivas de Formação Geral são padronizadas, ou seja, constam de todas as provas, os comentários da banca são os mesmos para todas as carreiras acadêmicas, sendo direcionados a todos os estudantes que participaram do ENADE/2010.
A seguir, serão analisados os desempenhos de Ingressantes e Concluintes da Área de Odontologia nas duas questões discursivas de Formação Geral do ENADE/2010. 3.3.1.1 Análise da Questão Discursiva 9 do Componente de Formação Geral Os dados obtidos a partir das respostas à questão 9 encontram -se na Tabela 3.13 e no Gráfico 3.10.
Nessa questão – de melhor desempenho dentre as duas de Formação Geral – os Ingressantes tiveram média 44,3, e os Concluintes, 49,3, resultando com diferença de 5,0 pontos entre os dois grupos.
Deixaram a questão em branco 20,9 \% dos Ingressantes e 15,2 \% dos Concluintes, e mesmo não tendo deixado -a em branco obtiveram nota zero, respectivamente, 1,4 \% e 1,1 \%.
Entre Ingressantes, a nota mediana foi 50,0 pontos, enquanto Concluintes obtiveram 55,0.
Estes últimos tiveram maior proporção de notas Não nulas no intervalo de 60 a 70, com 17,5 \% dos estudantes nessa faixa.
A nota máxima, 100,0 pontos, foi alcançada tanto pelos Concluintes quanto pelos Ingressantes. 3.3.1.2 Comentários sobre as respostas à Questão Discursiva 9 No que tange mais precisamente à atualidade do tema proposto e ao desenvolvimento do conteúdo, deve ser ressaltado que se mostrou sintonizado com a realidade política brasileira, bem como com a necessária reflexão sobre o papel na sociedade do profissional formado pelas Instituições de Ensino Superior do país.
Dissertar sobre democracia e ética foi algo apropriado para a época em que a prova foi aplicada, pois coincidiu com o período das eleições, em que tais temas são mais discutidos na vida cotidiana.
Não se pode, entretanto, deixar de se levar em conta a tendência de os estudantes terem limitado a associação do tema “ ética ” ao assunto “ política ”, considerando -se que a ética abarca não apenas a política.
Tal associação atingiu, em alguma medida, o alcance mais preciso das respostas.
As respostas apresentadas refletiram motivação para escrever sobre o tema da questão.
Alguns estudantes chegaram a desenvolver textos \textbf{eloquentes} e engajados.
Em alguns casos, o conceito de democracia, por exemplo, foi bem desenvolvido ao longo do texto, simultaneamente às exemplificações pedidas.
Além disso, muitos estudantes associaram o exercício da ética ao seu futuro ofício, reconhecendo o seu papel como cidadão no desempenho de sua profissão.
Essa associação insinuou, também, o esforço de reconhecimento da noção de ética prática, ainda que de maneira \textbf{codificada}.
Um dos pontos positivos revelou -se no fato de, na maioria das respostas analisadas, haver uma preocupação em definir a sociedade democrática.
Isso facilitou a correção e elevou a média de notas.
Outro aspecto positivo foi a presença pertinente de exposições coerentes de comportamento ético de um futuro profissional.
A articulação dos itens “ a ” e “ b ” da questão produziu, portanto, todo tipo de reflexão acerca de uma postura ética na sociedade, inclusive, corporativa.
A comprovação de que o tema ética tem sido abordado nos cursos superiores foi outro destaque positivo detectado durante as correções.
Assim, além do caráter crítico de que muitas das respostas corretas se revestiram, os relatos constantes sobre desdobramentos da ética comprovaram a abordagem sobre o tema, no mínimo, a ética profissional, nas carreiras universitárias que participaram do ENADE.
A despeito de alguns estudantes se posicionarem contra o exame, seja anulando a questão ou escrevendo algum comentário desabonador sobre o processo, foi observado que significativa parte dos participantes tem consciência político-acadêmica e percebe por que está realizando aquele tipo de prova.
Desse modo, a maioria dos estudantes tinha algo a “ dizer “ sobre o tema da questão, isto é, valia -se da oportunidade para se posicionar a respeito do assunto trazido à \textbf{baila}.
Ainda do ponto de vista do conteúdo desenvolvido nas respostas, observou -se preocupação solidária e consciência cidadã dos estudantes.
Mais um destaque positivo crucial se refere à habilidade textual.
Os estudantes, por exemplo, expressaram atitudes éticas e não éticas e as contrapunham ao conceito de democracia.
Há casos em que os alunos definiam a corrupção política como um exemplo de atitude não ética, enfatizando a real missão desse cargo em uma sociedade “ atender aos anseios da população e zelar pelo bem comum ”, escreveram vários estudantes.
No tocante à produção de texto dissertativo, à objetividade e ao caráter sintético, muitos estudantes corresponderam ao solicitado, elaborando pequenos textos dissertativos.
Nessas respostas, percebeu -se demonstração de domínio do assunto focado e utilização dos textos de referência como base para a elaboração.
O enunciado da questão 9 determinou que fosse elaborado um texto dissertativo que contivesse a )..., b )... e c )..., fornecendo um encaminhamento para as respostas que conduziu tanto a um resultado positivo quanto a um negativo, a saber : foi produtiva a resposta dos estudantes no sentido de que conseguiram, na maior parte das vezes, atender plenamente aos três quesitos apresentados na questão ; por outro lado, mais especificamente devido ao uso das três letras indicativas ( a, b e c ), houve tendência à elaboração de respostas divididas em itens, em vez de um texto eminentemente dissertativo.
Muitos estudantes apresentaram respostas nesse formato fragmentado para satisfazer a formalidade do roteiro sugerido no comando da questão.
A ideia-chave que deveria permear toda a resposta – democracia – se perdia em fragmentos desconexos entre si.
Se, por um lado, o roteiro do comando da questão facilitou o raciocínio de alguns dos estudantes, por outro, induziu grande parte deles a dois caminhos : i ) o das respostas curtas e objetivas, do tipo “ um profissional de saúde não ético é aquele que conta para as outras pessoas o que foi feito na consulta com o seu paciente ” e ii ) o outro caminho, que é do contraponto entre o conceito de sociedade democrática e a realidade brasileira.
Esse segundo caminho levou vários universitários a produzirem textos bem escritos, distantes, entretanto, do tema exigido.
No que se reporta às especificidades do desenvolvimento do conteúdo, a dificuldade ou falta de competência comunicativa revelada em grande parte das respostas denuncia a pouca exposição dos alunos à leitura de textos escritos de mais excelência.
Tal dificuldade mostrou -se mais evidente na impropriedade no uso de conectores e na incapacidade de elaboração de texto dissertativo, no tratamento de questões mais abstratas.
Muitos estudantes demonstraram não saber em que consistia a ética e confundiram exemplificação com definição : ao tentarem dar exemplos, via de regra, elaboravam uma definição incipiente do que se pedia.
Portanto, a dificuldade dos estudantes em argumentar e articular as ideias com coesão e coerência ainda são os grandes “ vilões ” desse processo avaliativo.
O excesso de fragmentos e textos prévios acerca do tema permitiu que a elaboração das respostas se resumisse a uma colagem ou paráfrase mal feita dos materiais disponíveis.
Em vez de subsidiar as respostas, esse material acabou por influenciar negativamente o desempenho dos estudantes.
Desse modo, verificou -se um nível de letramento ainda significativamente baixo, visto que a maioria apresentou uma visão bastante reducionista e simplificada sobre o tema em pauta.
De um modo geral, ao conceituarem ética e sociedade democrática, os estudantes demonstraram desconhecer o limite entre os dois conceitos, sem aludirem à interseção de ambos.
Foi expressiva a quantidade de redações nas quais questionavam a democracia pelo simples fato de estarem fazendo, a fortiori, aquela prova especificamente.
Em muitas das respostas, os estudantes se diziam incomodados pelo fato de estarem obrigados a realizar o exame, demonstrando sua insatisfação.
Por diversas vezes, os estudantes não foram capazes de sintetizar e de selecionar seu discurso o suficiente para atender ao que era pedido, o que ocasionou a falta de, pelo menos, um dos aspectos da resposta.
Foi constatada grande parte de respostas com problemas graves em termos textuais e gramaticais, em relação ao que se esperaria de pessoas que estão vinculadas ao Ensino ( Técnico ) Superior.
Alguns estudantes sequer conseguiam elaborar uma frase coerente e coesa, apresentando respostas que poderiam ser classificadas como de analfabetismo funcional.
De maneira ampla, manifestou -se da parte dos estudantes uma grande dificuldade no trato da língua escrita, dificuldade essa que se traduz, sobretudo, por um desrespeito bastante frequente às regras relativas : i ) à ortografia ( por exemplo : “ proficionais ”, “ infligir as leis ”, “ corrupitos ”, “ democrácia ”, “ opitar ”, “ manter o cizilo proficionau ” ), troca de “ sociedade democrática ” por “ sociedade demográfica ” e de “ ética ” por “ étnica ” ; ii ) à pontuação ; iii ) à concordância verbal ( “ os político deve (... ) eles pode ” ) ; iv ) ao modo de organização textual ( sem distinção entre o que seja dissertativo, argumentativo ou narrativo ).
Frisa -se que esse comentário não se restringe ao que se costuma chamar de “ língua padrão ”, mas ao que poderia ser chamado de “ excelência textual ”, que não se mede exclusiva ou enfaticamente pela correção gramatical, mas pela competência, objetividade, clareza e, em alguma medida, pela criatividade no ato de redigir.
Por fim, destaca -se um dos problemas dignos de nota na formulação da questão.
A transcrição do Dicionário Houaiss apresentava defeitos.
Por exemplo, a definição 5 de “ democracia ” indicava que esse é o tipo de “ governo que acata a vontade da maioria, embora respeitando os direitos e a livre expressão das minorias ”.
O emprego de um valor concessivo na oração adverbial, por oposição a um valor aditivo, contradiz as concepções efetivamente democráticas apresentadas de 1 a 4.
Repercussões desse valor concessivo, mal empregado, foram encontradas em incontáveis respostas, as quais diziam ser típico de um governo democrático, seguir a vontade da maioria, “ apesar ” de respeitar as minorias.
Outro desses problemas dignos de nota refere -se à indicação do 3.3.1.3 Análise da Questão Discursiva 10 do Componente de Formação Geral O desempenho dos estudantes na questão 10 foi inferior ao obtido na de número 9, com nota média 36,7 para Ingressantes e 41,7 para Concluintes.
A nota máxima também foi atingida em ambos os grupos, sendo que os percentuais de alunos que deixaram a questão em branco foram de 26,3 \% e 20,5 \%, respectivamente, para Ingressantes e Concluintes.
Tais índices revelam que houve maior tendência a deixar a questão 10 em branco, uma vez sendo comparada à de número 9.
Das notas não nulas, houve maior percentual de alunos, considerando -se a distribuição de notas, no intervalo entre 40 e 50 pontos, tanto para Ingressantes quanto para Concluintes : de 15,0 \% e 16,9 \%, respectivamente.
Não se manifestou uma superioridade acentuada de Concluintes nas diferentes faixas de notas posteriores a esse intervalo, com equilíbrio no desempenho dos dois grupos.
Essas informações estão ilustradas na Tabela 3.14 e no Gráfico 3.11. 3.3.1.4 Comentários sobre as respostas à Questão Discursiva 10 Inicialmente, será apresentada uma panorâmica sobre as respostas, com destaque para os aspectos positivos observados.
Em relação ao desenvolvimento do tema e à vinculação do mesmo à vida cotidiana dos estudantes que participaram do exame, a visão positiva em relação à Educação de Jovens e Adultos ( EJA ) reforça a ideia de que a educação pode trazer crescimento não somente profissional, mas, sobretudo, para ampliar os \textbf{horizontes} da cidadania.
Foram frequentes as respostas que apontavam a educação como a melhor saída, senão a única, para o resgate da cidadania.
A atuação dos governos, em especial o governo federal, foi um ponto destacado nas respostas, as quais sinalizavam para a necessidade de as instâncias governamentais assumirem um papel cada vez mais ativo na consolidação da educação – pública, sobretudo – no país.
As políticas públicas focadas na educação foram, então, um item bastante frequente nas respostas, seja para enaltecer, criticar ou requisitar a ação governamental.
Outro ponto positivo destacado diz respeito à consciência ou ao reconhecimento da precariedade de condições de trabalho dos professores, em geral, no Brasil e, em particular, daqueles ocupados com turmas de jovens e adultos.
Nesse sentido, foram inúmeras as propostas apresentadas na segunda parte das respostas com relação à imperiosidade de melhoria salarial para os docentes, e à melhoria das condições gerais de trabalho desses profissionais.
A associação direta entre valorização do professor e sucesso escolar dos estudantes foi um ponto complementar a essa preocupação com os docentes.
Os estudantes demonstraram uma visão de vínculo no processo de ensino-aprendizagem, no qual alunos e professores são partes inalienáveis entre si.
A questão da repetência e da evasão escolar, referenciadas no suporte da questão, foram então, abordadas pelos alunos por meio do viés dessa valorização docente.
Ainda que os temas específicos da repetência e da evasão não fossem contemplados a contento na maior parte das respostas, ao menos, nessa associação com a qualificação e a valorização dos professores, algo foi trabalhado analiticamente pelos estudantes.
Num plano mais abrangente, houve reconhecimento de que, no processo ensino-aprendizagem, não há atores isolados : alunos, profissionais de educação, governos e a própria sociedade como \textbf{corresponsáveis} pelo sucesso ou fracasso de tal processo.
A atualidade do tema e a sua explícita vinculação com o cotidiano de muitos estudantes que realizavam o exame fizeram com que os respondentes demonstrassem um visível grau de envolvimento com o tema proposto.
Ou seja, foi possível detectar um empenho notório dos estudantes ao redigir suas respostas com seriedade e comprometimento relativo à delicada problemática dos milhões de estudantes trabalhadores do país.
Utilizar todo o espaço destinado à resposta, procurando aproveitar ao máximo a oportunidade de emitir um argumento ou uma opinião revelou -se como uma das evidências do aludido empenho.
A EJA também foi reconhecida como um dos caminhos concretos para o acesso à Universidade, no médio prazo.
As respostas se voltaram, muitas vezes, à condição de depoimentos pessoais sobre a conciliação de trabalho e estudo.
Nessa direção, ao assumirem o papel de depoentes, os estudantes expuseram suas histórias de vida, suas expectativas e percalços, suas \textbf{esperanças} e \textbf{entraves} cotidianos, seus anseios mais ambiciosos e problemas estruturais, que poderiam se resolver por meio da educação.
Nesse inciso, ficava clara a relação que se formava entre educação e mobilidade social, mediatizada por iniciativas governamentais como a EJA.
Uma sugestão plausível e válida para a solução dos problemas com a EJA no país reportou -se à necessidade de maior participação do público alvo na definição de estratégias para essa modalidade de ensino.
Nesse detalhe, era notória a associação das preocupações relativas à questão 9, que trata do problema da democracia e da cidadania, com o problema específico da educação escolar.
Houve um número significativo de sugestões ou propostas para se combater o problema da repetência e da evasão, destacando -se melhores salários para os professores e aulas mais dinâmicas.
A lista dessas sugestões, contudo, era bem mais ampla e diversificada, incluindo ainda : a redução da jornada de trabalho dos jovens e adultos que participam da EJA ; melhoria de infraestrutura nas instituições responsáveis pela EJA ; incentivo financeiro aos estudantes trabalhadores por meio de um sistema de bolsas de estudo ; realização de campanhas motivacionais para o retorno aos estudos ; estímulo à participação das empresas, por meio da solicitação de recursos materiais, espaço físico para as salas, auxílio financeiro ; incremento dos cursos profissionalizantes no país ; e ampliação da oferta de ensino a distância em todo o território nacional.
Muitos estudantes abordaram, com propriedade, a associação entre um mercado de trabalho cada vez mais competitivo e a necessidade de qualificação profissional, inclusive com uma clara percepção de que os salários são tanto maiores quanto maior a quantidade de anos de estudo e/ou a qualidade da formação.
Nesse tópico, retorna -se ao papel crucial da EJA como viabilizadora dessa formação.
Em continuidade à análise, serão apresentados aspectos negativos que foram notados com frequência nas respostas à questão 10.
De um ponto de vista mais global, a dificuldade repousa no entendimento da questão, isto é, na decifração do comando.
Muitos fugiram completamente do assunto proposto, ainda que o enunciado da questão fosse farto em dados específicos ao tema.
Em vez de responderem diretamente quais eram os problemas enfrentados por aqueles que têm de trabalhar e estudar ao mesmo tempo, muitos estudantes optavam por falar das dificuldades que existem na educação brasileira, de forma estendida e pasteurizada.
Vários estudantes não conseguiam, em suas respostas, relacionar as dificuldades de conciliação da escola com o trabalho, a não ser, como já mencionado, nos casos de depoimentos pessoais, relatos de experiência, ou seja, traços de uma biografia.
Assim, eram respostas cujo teor era apenas o de registrar o abandono dos estudos por parte de vários jovens que têm de trabalhar.
O tom das respostas, nesses casos, era meramente de uma descrição e não de uma análise, muito menos de um elemento que compusesse uma argumentação.
A dificuldade de compreensão da questão se expressava nas respostas que, em lugar de mostrarem os problemas de conciliação entre escola e trabalho, prendiam -se à ideia de “ associação ” – termo que consta do enunciado da questão – entre escola e trabalho.
Muitos alunos relatavam as consequências positivas de uma boa educação no momento em que se busca bons postos de trabalho.
Tratava -se de uma fuga parcial ao comando da questão, evidenciando má compreensão do mesmo.
Uma dificuldade facilmente percebida nas respostas refere -se à tendência de os estudantes tratarem dos problemas relativos à escola e, muito separadamente, daqueles problemas relativos ao trabalho, reconhecendo uma espécie de sinergia que impediria de visualizá -los em sua individualidade, ao mesmo tempo que em sua complementaridade.
Isso significa que os estudantes distinguiam por mero isolamento, ou seja, eles não ampla e diversificada, incluindo ainda : a redução da jornada de trabalho dos jovens e adultos que participam da EJA ; melhoria de infraestrutura nas instituições responsáveis pela EJA ; incentivo financeiro aos estudantes trabalhadores por meio de um sistema de bolsas de estudo ; realização de campanhas motivacionais para o retorno aos estudos ; estímulo à participação das empresas, por meio da solicitação de recursos materiais, espaço físico para as salas, auxílio financeiro ; incremento dos cursos profissionalizantes no país ; e ampliação da oferta de ensino a distância em todo o território nacional.
Muitos estudantes abordaram, com propriedade, a associação entre um mercado de trabalho cada vez mais competitivo e a necessidade de qualificação profissional, inclusive com uma clara percepção de que os salários são tanto maiores quanto maior a quantidade de anos de estudo e/ou a qualidade da formação.
Nesse tópico, retorna -se ao papel crucial da EJA como viabilizadora dessa formação.
Em continuidade à análise, serão apresentados aspectos negativos que foram notados com frequência nas respostas à questão 10.
De um ponto de vista mais global, a dificuldade repousa no entendimento da questão, isto é, na decifração do comando.
Muitos fugiram completamente do assunto proposto, ainda que o enunciado da questão fosse farto em dados específicos ao tema.
Em vez de responderem diretamente quais eram os problemas enfrentados por aqueles que têm de trabalhar e estudar ao mesmo tempo, muitos estudantes optavam por falar das dificuldades que existem na educação brasileira, de forma estendida e pasteurizada.
Vários estudantes não conseguiam, em suas respostas, relacionar as dificuldades de conciliação da escola com o trabalho, a não ser, como já mencionado, nos casos de depoimentos pessoais, relatos de experiência, ou seja, traços de uma biografia.
Assim, eram respostas cujo teor era apenas o de registrar o abandono dos estudos por parte de vários jovens que têm de trabalhar.
O tom das respostas, nesses casos, era meramente de uma descrição e não de uma análise, muito menos de um elemento que compusesse uma argumentação.
A dificuldade de compreensão da questão se expressava nas respostas que, em lugar de mostrarem os problemas de conciliação entre escola e trabalho, prendiam -se à ideia de “ associação ” – termo que consta do enunciado da questão – entre escola e trabalho.
Muitos alunos relatavam as consequências positivas de uma boa educação no momento em que se busca bons postos de trabalho.
Tratava -se de uma fuga parcial ao comando da questão, evidenciando má compreensão do mesmo.
Uma dificuldade facilmente percebida nas respostas refere -se à tendência de os estudantes tratarem dos problemas relativos à escola e, muito separadamente, daqueles problemas relativos ao trabalho, reconhecendo uma espécie de sinergia que impediria de visualizá -los em sua individualidade, ao mesmo tempo que em sua complementaridade.
Isso significa que os estudantes distinguiam por mero isolamento, ou seja, eles não conseguiam distinguir e articular e, sobretudo, reconhecer as convocações mútuas entre os problemas da escola e do trabalho, no caso explícito dos estudantes trabalhadores.
Essa dificuldade de identificar, associar e analisar aspectos distintos, até mesmo concorrentes, antagônicos e complementares, denuncia a fragilidade de um pensamento simplificador que a qualidade das respostas revelou ser ainda hegemônico nas universidades brasileiras.
Notou -se a presença de preconceitos na análise dos fatores que levam ao abandono dos estudos.
Muitas respostas apontavam como causa desse abandono e da repetência a preguiça, a falta de perseverança e o cansaço dos estudantes.
Nesse tópico específico, fica evidente a responsabilização do indivíduo pelos percalços e fracassos de sua vida, descontextualizando -o de uma sociedade que, em seu movimento de conjunto, lhe impõe condicionantes.
A precariedade do uso da língua pátria de modo culto foi uma constante, conforme mencionado na análise da questão 9.
As impropriedades lexicais e a falta de elementos coesores que articulassem as partes constitutivas das respostas foram um incômodo constante.
Nesse sentido, agrava -se a situação da baixa qualidade textual em decorrência complementar da construção de textos meramente descritivos e não argumentativos, como é exigido no suporte da questão.
Isso demonstra a dificuldade dos estudantes universitários na compreensão de que uma argumentação requer, imperiosamente, justificações para os apontamentos apresentados.
Os textos das respostas eram, em sua maioria, opinativos, sem qualquer compromisso com elementos explicativos ou com a busca de justificativas. 3.3.2 Componente de Conhecimento Específico Na parte da prova relativa ao componente de Conhecimento Específico, observa -se que a diferença de desempenho foi mais acentuada entre Ingressantes e Concluintes.
A Tabela 3.15 contém dados elucidativos, na qual é possível observar que, no grupo dos Ingressantes, a nota média 6,8 é bem inferior à alcançada por Concluintes, 40,9.
Essas médias foram menores do que as obtidas no conjunto das questões objetivas, nas quais Ingressantes alcançaram nota média 37,2, e Concluintes, 57,9 ( Tabela 3.8 ).
O Gráfico 3.12 apresenta a distribuição de notas nos diferentes intervalos.
Observa -se que os Concluintes atingiram proporções mais elevadas de notas na faixa entre 40 e 50 pontos, com 21,6 \% dos estudantes.
A maior concentração de Ingressantes é identificada no intervalo entre zero e 10 pontos, correspondendo a 74,7 \% desse grupo.
A análise de cada uma destas questões será feita a seguir. ∗ 3.3.2.1 Análise da Questão Discursiva 38 do Componente de Conhecimento Específico Na questão 38, cujos resultados aferidos encontram -se descritos na Tabela 3.16, o grupo dos Ingressantes obteve nota média 2,8, enquanto a média dos Concluintes foi ∗ A análise completa, com apresentação de trechos de respostas de estudantes, pode ser encontrada no documento “ Relatório da Correção das Questões Discursivas ”, elaborado pelos docentes que participaram do processo de correção das mesmas. grupo dos Ingressantes obteve nota média 2,8, enquanto a média dos Concluintes foi 32,0.
Deixaram a questão em branco 65,2 \% dos Ingressantes e 15,5 \% dos Concluintes.
As proporções de notas dos Concluintes foram superiores às dos Ingressantes a partir do intervalo de zero a 10 pontos, sendo a faixa modal aquela entre 20 e 30 pontos, com 14,5 \%.
A nota máxima, 100,0, foi alcançada apenas no grupo dos Concluintes, sendo 90,0 a máxima entre Ingressantes. 3.3.2.2 Comentários sobre as respostas à Questão Discursiva 38 Esta questão estava bem formulada, dividida em subitens, possibilitando ao aluno que dominava o assunto obter o grau máximo.
Apesar desses indicadores que favoreciam ao estudante, o número de respostas em branco foi bastante alto.
Quando correta, a resposta ao item “ a ” apresentou, na maioria das vezes, somente uma das duas hipóteses de diagnóstico denominadas de forma adequada à literatura e ao gabarito.
Exemplo 1 : “ Mucocele ou hiperplasia fibrosa ” ( 1.0 ponto ) ; Exemplo 2 : “ As duas hipóteses de diagnóstico podem ser fibroma de irritação ou um fibroma reacional ” ( 1.0 ponto ) ; Exemplo 3 : “ granuloma ou fibroma traumático ” ( 1.0 ponto ).
Muitas respostas apresentaram hipóteses compatíveis com câncer, completamente fora do gabarito, correspondente a zero na obtenção de pontos.
Considerando o item “ b ” ( justificativa para a necessidade ou não de descontinuar o uso da medicação previamente à biópsia ) as respostas apresentaram -se, em sua maioria, confusas e contraditórias.
Nenhuma questão corrigida atendeu ao gabarito completo, correspondente à resposta.
Muitos estudantes erraram a questão ao optarem por suspender o uso do medicamento.
Em relação ao item “ c ”, observou -se que a maioria dos acertos incluiu os dois itens ou “ sutura + compressão ” adequadamente respondidos.
Muitas respostas incluíram pelo menos a indicação de “ sutura ”.
Não foi constatada em nenhuma prova a resposta que envolvia o uso de enxaguatórios bucais com ácido épsilon-aminocapróico ou ácido tranexânico, proposta no gabarito.
Considerando o item “ d ”, muitas respostas incluíram pacientes com hipertensão, mas constantemente citavam problemas cardíacos, embolia pulmonar, válvulas cardíacas protéticas, uso de stents, AVE e trombose, com a obtenção de várias respostas completamente corretas. 3.3.2.3 Análise da Questão Discursiva 39 do Componente de Conhecimento Específico A Tabela 3.17 contém as informações relativas à questão 39 do conjunto de questões do componente de Conhecimento Específico.
O desempenho dos dois grupos – Ingressantes e Concluintes – nessa questão, foi superior ao desempenho na questão de número 38.
As médias foram 10,8 no grupo dos Ingressantes e 49,8 no grupo dos Concluintes.
A proporção de estudantes que deixaram a questão em branco foi menor do que a observada na questão 38, embora bastante alta, especialmente entre Ingressantes.
Estes corresponderam a 59,2 \%, enquanto entre Concluintes, 7,1 \% foram os que não responderam.
A nota máxima ( 100,0 ) foi atingida somente por Concluintes.
Entre Ingressantes, foi 85,0.
O intervalo de 20 a 30 pontos correspondeu ao intervalo de notas no qual houve maior equilíbrio entre os dois grupos : Ingressantes com 6,5 \% e Concluintes com 6,1 \%.
Nos intervalos seguintes, as notas de Ingressantes se mantiveram baixas, havendo crescimento da proporção de Concluintes atingindo a moda na faixa entre 50 e 60 pontos : 22,4 \% desse grupo.
As distribuições das notas entre os estudantes constam do Gráfico 3.14. 3.3.2.4 Comentários sobre as respostas à Questão Discursiva 39 A questão solicitava um planejamento de tratamento para o caso clínico apresentado.
As respostas em branco foram numerosas.
Apesar de contemplarem o gabarito, muitas respostas corretas apresentavam trechos errados que mostravam a falta de conhecimento do aluno em determinados aspectos.
Muitos estudantes, apesar de apresentarem um planejamento coerente, não responderam a todos os quesitos necessários para o grau máximo.
Algumas respostas foram extremamente simplistas para tratar o paciente descrito.
Ao reabilitar o paciente com bruxismo, sem encerramento diagnóstico, a placa miorrelaxante, por exemplo, pode acarretar o agravamento da condição.
No entanto, ao mencionar partes do gabarito divulgado, a resposta foi considerada no que foi possível.
Diversos estudantes dividiram o planejamento por especialidades.
Outro fato que chamou bastante atenção foi que muitos indicaram Prótese Fixa ( PF ) para repor os molares inferiores, o que está completamente errado.
De acordo com o exposto, nessa questão, quase nenhum estudante foi capaz de planejar corretamente uma sequência de tratamento reabilitador para o paciente, em conformidade com o gabarito fornecido. 3.3.2.5 Análise da Questão Discursiva 40 do Componente de Conhecimento Específico A questão de número 40 foi anulada pela Comissão constituída pelo INEP, motivo pelo qual não há dados de análise a serem expostos.
CONSIDERAÇÕES FINAIS Ao final da correção, os professores concluíram, a partir das respostas às duas questões disponíveis, que a capacidade de integrar diferentes áreas da Odontologia para que o paciente possa ser tratado de forma integral ainda é uma lacuna dos estudantes.
A Banca de Correção lamentou a anulação da questão 40.
Cabe salientar, ainda, que segundo avaliação da Banca, o grau de dificuldade das questões elaboradas para nível de graduação foi elevado.
Isso foi evidenciado pelas notas baixas apresentadas e pela quantidade de respostas em branco.
No entanto, os acordos firmados entre os professores corretores para uma melhor correção não deixaram nenhum estudante ou Instituição de ensino superior em prejuízo.
Entende -se, por fim, que um currículo que integre as disciplinas na prática clínica, parece ser urgentemente necessário às instituições de ensino de Odontologia. apresentarem um planejamento coerente, não responderam a todos os quesitos necessários para o grau máximo.
Algumas respostas foram extremamente simplistas para tratar o paciente descrito.
Ao reabilitar o paciente com bruxismo, sem encerramento diagnóstico, a placa miorrelaxante, por exemplo, pode acarretar o agravamento da condição.
No entanto, ao mencionar partes do gabarito divulgado, a resposta foi considerada no que foi possível.
Diversos estudantes dividiram o planejamento por especialidades.
Outro fato que chamou bastante atenção foi que muitos indicaram Prótese Fixa ( PF ) para repor os molares inferiores, o que está completamente errado.
De acordo com o exposto, nessa questão, quase nenhum estudante foi capaz de planejar corretamente uma sequência de tratamento reabilitador para o paciente, em conformidade com o gabarito fornecido. 3.3.2.5 Análise da Questão Discursiva 40 do Componente de Conhecimento Específico A questão de número 40 foi anulada pela Comissão constituída pelo INEP, motivo pelo qual não há dados de análise a serem expostos.
CONSIDERAÇÕES FINAIS Ao final da correção, os professores concluíram, a partir das respostas às duas questões disponíveis, que a capacidade de integrar diferentes áreas da Odontologia para que o paciente possa ser tratado de forma integral ainda é uma lacuna dos estudantes.
A Banca de Correção lamentou a anulação da questão 40.
Cabe salientar, ainda, que segundo avaliação da Banca, o grau de dificuldade das questões elaboradas para nível de graduação foi elevado.
Isso foi evidenciado pelas notas baixas apresentadas e pela quantidade de respostas em branco.
No entanto, os acordos firmados entre os professores corretores para uma melhor correção não deixaram nenhum estudante ou Instituição de ensino superior em prejuízo.
Entende -se, por fim, que um currículo que integre as disciplinas na prática clínica, parece ser urgentemente necessário às instituições de ensino de Odontologia.

% matched lemmas: baila, codificar, corresponsável, eloquente, entrave, esperança, horizonte
\end{textsample}
